\documentclass[12pt]{article}
\usepackage[utf8]{inputenc}
\usepackage[T1]{fontenc}
\usepackage{hyperref}
\title{Dumo Language Resources for K-12 Curriculum Development}
\author{Compiled by Claude Sonnet 5 (claude-sonnet-5) for Andr\'as Kornai}
\date{}
\begin{document}
\maketitle

\section{Introduction}
Dumo (ISO 639-3 vam; also called Vanimo, Manimo or Duso in some sources) is a tonal Skou language of Sandaun Province, Papua New Guinea. Sources found give estimates of roughly 2,000 to 2,700 speakers. Tok Pisin is reported as a secondary language. Searches found no school materials, primers or textbooks in Dumo. The resources below are reference and background only.

\section{Primary Educational Resources}
No primers, textbooks or school programmes in Dumo were found in the searches.

\section{Linguistic and Reference Materials}
\begin{itemize}
\item Wikipedia, Vanimo language (covers Dumo): \url{https://en.wikipedia.org/wiki/Vanimo_language}
\item Wikipedia, Dumo people: \url{https://en.wikipedia.org/wiki/Dumo_people}
\item Dumo numerals page, Max Planck Institute (EVA): \url{https://lingweb.eva.mpg.de/channumerals/Dumo.htm}
\item WALS entry for Dumo: \url{https://wals.info/languoid/lect/wals_code_dum}
\item Glottolog reference record for Malcolm Ross 1980 (Skou classification): \url{https://glottolog.org/resource/reference/id/72393}
\end{itemize}

\section{Religious and Community Resources}
\begin{itemize}
\item Joshua Project profile, Vanimo in Papua New Guinea: \url{https://joshuaproject.net/people_groups/15771/pp}
\end{itemize}
No Dumo scripture translation was identified in the search results.

\section{Digital and Audio Resources}
\begin{itemize}
\item Endangered Languages Project entry (lists the language as vulnerable): \url{https://endangeredlanguages.com/lang/10588?hl=es}
\end{itemize}
No audio or video collections were found in the searches.

\section{Curriculum Development Context}
Little usable material exists. A curriculum would need new work in orthography development, documentation and materials writing, with community involvement. The tonal system and the classification debate noted in the sources should be checked with specialists. Dumo is also associated with the name Vanimo, which is also the name of a town and of other groups, so resources should be checked against this language.
\end{document}